% Unit 070 terjemahan Bahasa Indonesia dari chapter5.tex baris 955--1063.
% Sumber: Methods of Algebra, Volume 2, commit
% 9a5803ff77dd3257484cb177f851a73770a59dd3 (CC BY 4.0).
% stable-unit-id: o014.aljabr2.chapter5.multiplicative-structures
% Cakupan: Bagian 5.7 lengkap, sebelum latihan.

% segment-id: o014.li2.chapter5.multiplicative-structures.h001
\section{Sekilas tentang struktur perkalian}\label{sec:multiplicative-SS}
\hypertarget{o014.aljabr2.chapter5.multiplicative-structures}{ }

% segment-id: o014.li2.chapter5.multiplicative-structures.p001
Agar pembahasan sederhana dan konkret, dalam bagian ini kita mengambil
gelanggang komutatif $\Bbbk$ dan $\mathcal{A}=\Bbbk\dcate{Mod}$; ini sudah
memadai bagi penerapan klasik. Modul-$\Bbbk$ kita sebut singkat sebagai modul,
aljabar-$\Bbbk$ sebagai aljabar, dan kita tulis
$\otimes:=\otimes_{\Bbbk}$.

% segment-id: o014.li2.chapter5.multiplicative-structures.p002
Misalkan $I$ monoid komutatif, dengan operasi biner ditulis secara aditif.
Definisi dari \cite[\S 7.4]{Li1} dapat diringkas sebagai berikut:
\begin{itemize}
	\item modul bergradasi-$I$ adalah modul dengan dekomposisi jumlah langsung
	$M=\bigoplus_{i\in I}M^i$; unsur-unsur dalam $M^i$ disebut unsur homogen
	berderajat $i$;
	\item aljabar bergradasi-$I$ adalah modul bergradasi yang dilengkapi
	homomorfisme $\mu:A\otimes A\to A$, yang juga ditulis sebagai perkalian
	$xy:=\mu(x\otimes y)$, sedemikian sehingga $A$ menjadi gelanggang dan
	\[ 1_A\in A^0,\qquad A^i\cdot A^j\subset A^{i+j},
	\qquad i,j\in I; \]
	jadi perkalian $\mu$ sepenuhnya ditentukan oleh data
	$\mu^{i,j}:A^i\otimes A^j\to A^{i+j}$;
	\item homomorfisme dan isomorfisme antara modul atau aljabar
	bergradasi-$I$ didefinisikan dengan cara biasa; subaljabar dan ideal
	bergradasi-$I$ juga didefinisikan dengan cara yang sama.
\end{itemize}

% segment-id: o014.li2.chapter5.multiplicative-structures.p003
Untuk $k\in I$, definisikan pula homomorfisme berderajat $k$ antara modul
bergradasi-$I$: ini adalah homomorfisme modul $\varphi:M\to N$ yang memenuhi
$\varphi(M^i)\subset M^{i+k}$ untuk semua $i$, dan ditentukan oleh data
$(\varphi^i:M^i\to M^{i+k})_{i\in I}$. Untuk $k=0$ kita memperoleh kembali
homomorfisme dalam arti biasa. Derajat homomorfisme dijumlahkan ketika
homomorfisme dikomposisikan.

% segment-id: o014.li2.chapter5.multiplicative-structures.d001
\begin{definition}\label{def:dg-algebra-preview}
	\index{aljabar bergradasi diferensial (differential graded algebra)}
	Misalkan $I$ monoid komutatif yang dilengkapi homomorfisme
	$\epsilon:I\to\Z/2\Z$. Sebuah \emph{aljabar bergradasi-$I$ diferensial}
	dengan diferensial berderajat $k$ adalah aljabar bergradasi-$I$
	$A=\bigoplus_{p\in I}A^p$ bersama endomorfisme berderajat $k\in I$
	$d=(d^p)_p:A\to A$ yang memenuhi $d^2=0$ dan aturan Leibniz
	\[ d(xy)=(dx)\cdot y+(-1)^{\epsilon(p)}x\cdot dy,
	\qquad x\in A^p,\quad y\in A^{p'},\quad p,p'\in I. \]
\end{definition}

% segment-id: o014.li2.chapter5.multiplicative-structures.p004
Perhatikan bahwa
$d(1_A)=d(1_A\cdot1_A)=d(1_A)+d(1_A)$, sehingga $d(1_A)=0$. Jelas bahwa
$\Ker(d)$ merupakan subaljabar bergradasi-$I$ dari $A$, sedangkan
$\Image(d)$ merupakan ideal dua sisi bergradasi-$I$ dalam $\Ker(d)$. Karena
itu,
\[ \Hm(A,d):=\Ker(d)/\Image(d) \]
tetap merupakan aljabar bergradasi-$I$.

% segment-id: o014.li2.chapter5.multiplicative-structures.e001
\begin{example}
	Ambil $I=\Z$ dan $\epsilon(p)=p\;\bmod 2$. Untuk setiap kompleks $X$ pada
	kategori $\Bbbk$-linear $\mathcal{B}$, kompleks $\Hom$
	$\Hom^\bullet(X,X)$ dari Definisi~\ref{def:Hom-cplx}, bersama
	$d=d_{\Hom^\bullet(X,X)}$, menjadi aljabar bergradasi-$\Z$ diferensial
	dengan diferensial berderajat $1$. Pernyataan ini pada dasarnya merupakan
	isi Lema~\ref{prop:Hom-cplx-Leibniz}. Kita mempunyai
	\[ \Hm(\Hom^\bullet(X,X),d)=
	\bigoplus_{n\in\Z}\Hom_{\cate{K}(\mathcal{B})}(X,X[n]), \]
	dengan perkalian yang berasal dari komposisi morfisme dalam
	$\cate{K}(\mathcal{B})$.
\end{example}

% segment-id: o014.li2.chapter5.multiplicative-structures.e002
\begin{example}
	Ambil $\Bbbk=\CC$, $I=\Z$, dan $\epsilon(p)=p\;\bmod2$. Bentuk
	diferensial bernilai $\CC$ berderajat $p$ pada manifold mulus
	$\mathfrak{X}$ membentuk ruang vektor $A^p(\mathfrak{X})$. Definisikan
	$A(\mathfrak{X}):=\bigoplus_{p=0}^{\dim X}A^p(\mathfrak{X})$. Dengan
	perkalian $\mu(\omega\otimes\eta):=\omega\wedge\eta$ dan diferensial luar
	$\dd$, objek ini menjadi aljabar bergradasi-$\Z$ diferensial dengan
	diferensial berderajat $1$. Aljabar terkait
	$\Hm(A(\mathfrak{X}),\dd)=\bigoplus_p\Hm^p_{\mathrm{dR}}(\mathfrak{X})$
	tepat merupakan kohomologi de Rham dari $\mathfrak{X}$; perkalian yang
	berasal dari $\wedge$ memberinya struktur aljabar bergradasi. Ini merupakan
	objek fundamental dalam topologi dan isomorfik sebagai aljabar bergradasi
	dengan cincin kohomologi singular $\mathfrak{X}$.

	Perkalian pada $\Hm(A(\mathfrak{X}),\dd)$ juga memenuhi
	$xy=(-1)^{pq}yx$, dengan $x$ dan $y$ masing-masing unsur homogen berderajat
	$p$ dan $q$, karena perkalian pada $A(\mathfrak{X})$ sudah mempunyai sifat
	ini. Sebagai contoh sederhana yang indah, ruang proyektif kompleks
	$\mathfrak{X}:=\mathbb{P}^n(\CC)$ menghasilkan aljabar bergradasi
	$\Hm(A(\mathfrak{X}),\dd)\simeq\CC[t]/(t^{n+1})$, dengan peubah $t$
	bersesuaian dengan unsur homogen berderajat $2$.
\end{example}

% segment-id: o014.li2.chapter5.multiplicative-structures.p005
Kembali ke pokok pembahasan, dalam bagian ini kita mempertimbangkan kasus-kasus
berikut:
\begin{center}\begin{tabular}{|c|c|c|} \hline
	Singkatan & $I$ & $\epsilon:I\to\Z/2\Z$ \\ \hline
	Bergradasi tunggal & $\Z$ & $\epsilon(p)=p\;\bmod2$ \\
	Bergradasi ganda & $\Z^2$ & $\epsilon(p,q)=p+q\;\bmod2$ \\ \hline
\end{tabular}\end{center}

% segment-id: o014.li2.chapter5.multiplicative-structures.p006
Teori ini beririsan dengan Definisi~\ref{def:differential-object},
tetapi fokus bagian ini adalah perkalian yang belum dibahas sebelumnya.

% segment-id: o014.li2.chapter5.multiplicative-structures.p007
Untuk meringankan notasi, selanjutnya berbagai kasus tersebut secara
bersama-sama kita sebut ``bergradasi''; pembeda utamanya adalah
superskrip $p$ atau $(p,q)$. Demikian pula, aljabar bergradasi-$I$
diferensial secara bersama-sama disebut aljabar bergradasi diferensial.

% segment-id: o014.li2.chapter5.multiplicative-structures.c001
\begin{convention}
	Untuk barisan spektral bergradasi ganda kohomologis $\mathscr{E}$ dalam
	kasus $\mathcal{A}=\Bbbk\dcate{Mod}$, kita mengidentifikasi setiap halaman
	$E_r$ dengan modul bergradasi $\bigoplus_{p,q}E_r^{p,q}$ dan memandang
	$d_r=(d_r^{p,q})_{p,q}$ sebagai endomorfisme berderajat $(r,-r+1)$ pada
	modul bergradasi $E_r$. Kasus homologis dan kasus bergradasi tunggal
	diperlakukan dengan cara serupa.
\end{convention}

% segment-id: o014.li2.chapter5.multiplicative-structures.p008
Singkatnya, barisan spektral dengan struktur perkalian adalah barisan spektral
yang terdiri atas aljabar bergradasi diferensial.

% segment-id: o014.li2.chapter5.multiplicative-structures.d002
\begin{definition}\label{def:mult-ss}
	\index{barisan spektral!struktur perkalian (multiplicative structure)}
	Sebuah \emph{struktur perkalian} pada barisan spektral bergradasi ganda
	kohomologis $\mathscr{E}$ adalah keluarga homomorfisme
	$\mu_r:E_r\otimes E_r\to E_r$ (yakni ``perkalian'') sedemikian sehingga
	\begin{itemize}
		\item setiap $(E_r,\mu_r,d_r)$ menjadi aljabar bergradasi diferensial,
		dengan diferensial berderajat $(r,-r+1)$;
		\item $t_{r+1}:\Hm(E_r,d_r)\rightiso E_{r+1}$ dalam data barisan
		spektral merupakan isomorfisme aljabar bergradasi.
	\end{itemize}
	Struktur perkalian pada barisan spektral bergradasi tunggal atau homologis
	didefinisikan dengan cara serupa.
\end{definition}

% segment-id: o014.li2.chapter5.multiplicative-structures.p009
Pembahasan sebelumnya menunjukkan bahwa $Z_r=\Ker(d_r)$ merupakan subaljabar
bergradasi dari $E_r$, sedangkan $B_r=\Image(d_r)$ merupakan ideal dua sisi
bergradasi dari $Z_r$. Karena itu $\Hm(E_r,d_r)$ menjadi aljabar bergradasi;
	inilah makna tepat Definisi~\ref{def:mult-ss}.

% segment-id: o014.li2.chapter5.multiplicative-structures.p010
Lebih lanjut, jika limitnya ada, maka
$Z_\infty=\bigoplus_{p,q}Z_\infty^{p,q}$ juga merupakan subaljabar bergradasi
dengan $B_\infty=\bigoplus_{p,q}B_\infty^{p,q}$ sebagai ideal bergradasinya.
Jadi $E_\infty$ juga merupakan aljabar bergradasi.

% segment-id: o014.li2.chapter5.multiplicative-structures.p011
\index{aljabar-dg}
Sebagai contoh, perhatikan aljabar bergradasi diferensial yang diferensialnya
berderajat $1$, yang disingkat sebagai aljabar-dg. Dengan mengembangkan
definisi, ini setara dengan kompleks modul-$\Bbbk$
$X=(X^n,d_X^n)_{n\in\Z}$ sedemikian sehingga
$X\xlongequal{\text{diidentifikasi}}\bigoplus_nX^n$ dilengkapi perkalian
$\mu:X\otimes X\to X$ yang memenuhi $X^p\cdot X^{p'}\subset X^{p+p'}$ dan
aturan Leibniz
\[ d(xy)=dx\cdot y+(-1)^px\cdot dy,
	\qquad x\in X^p,\quad y\in X^{p'}. \]
Sekarang andaikan data $(X,d)$ diperluas menjadi kompleks terfiltrasi
$(X,d,\mathrm{F}^\bullet X)$. Ingat bahwa ini sudah berarti
$d(\mathrm{F}^pX)\subset\mathrm{F}^pX$; kita syaratkan pula agar perkalian
pada $X$ kompatibel dengan filtrasi:
\[ \mathrm{F}^pX^n\cdot\mathrm{F}^{p'}X^{n'}
	\subset\mathrm{F}^{p+p'}X^{n+n'}. \]
Dalam keadaan ini $(X,d,\mu,\mathrm{F}^\bullet X)$ disebut aljabar bergradasi
diferensial terfiltrasi. Untuk filtrasi terinduksi pada $\Hm(X,d)$, syarat di
atas memastikan bahwa
\[ \gr\Hm(X,d)\xlongequal{\text{diidentifikasi}}
	\bigoplus_{(p,q)\in\Z^2}\gr^p\Hm^{p+q}(X,d) \]
secara alami menjadi aljabar bergradasi.
\index{aljabar bergradasi diferensial!terfiltrasi (filtered)}

% segment-id: o014.li2.chapter5.multiplicative-structures.t001
\begin{proposition}\label{prop:dga-mult-ss}
	Misalkan $(X,d,\mu,\mathrm{F}^\bullet X)$ adalah aljabar bergradasi
	diferensial terfiltrasi, dengan diferensial berderajat $1$.
	\begin{enumerate}[(i)]
		\item Barisan spektral $\mathscr{E}$ dari kompleks terfiltrasi
		$(X,d,\mathrm{F}^\bullet X)$ mempunyai struktur perkalian kanonik.
		\item Isomorfisme
		$E_\infty^{p,q}\simeq\gr^p\Hm^{p+q}(X,d)$ dalam Teorema Konvergensi
		Klasik~\ref{prop:classical-convergence} sesungguhnya memberikan
		isomorfisme aljabar bergradasi
		$E_\infty\simeq\gr\Hm(X,d)$.
	\end{enumerate}
\end{proposition}
\begin{proof}
	Ambil $x\in E_r^{p,q}$ dan $y\in E_r^{p',q'}$. Berdasarkan deskripsi
	Proposisi~\ref{prop:filtered-cplx-ss}, pilih prapeta
	\begin{gather*}
		\widetilde{x}\in\mathrm{F}^pX^{p+q}\cap
		d^{-1}\left(\mathrm{F}^{p+r}X^{p+q+1}\right),\qquad
		\widetilde{y}\in\mathrm{F}^{p'}X^{p'+q'}\cap
		d^{-1}\left(\mathrm{F}^{p'+r}X^{p'+q'+1}\right).
	\end{gather*}
	Menurut definisi,
	\begin{gather*}
		\widetilde{x}\widetilde{y}\in
		\mathrm{F}^{p+p'}X^{p+q+p'+q'},\\
		d(\widetilde{x}\widetilde{y})
		=d\widetilde{x}\cdot\widetilde{y}
		+(-1)^{p+q}\widetilde{x}\cdot d\widetilde{y}
		\in\mathrm{F}^{p+p'+r}X^{p+q+p'+q'+1}.
	\end{gather*}
	Jadi $\widetilde{x}\widetilde{y}$ menentukan unsur
	$xy\in E_r^{p+p',q+q'}$. Perhitungan rutin (silakan periksa) menunjukkan
	bahwa $xy$ hanya bergantung pada $x$ dan $y$.

	Selanjutnya, $d_r:E_r\xrightarrow{(r,-r+1)}E_r$ diinduksi oleh diferensial
	$d$ pada $X$. Ini memberikan struktur aljabar bergradasi diferensial pada
	$E_r$; sifat asosiatif, aturan Leibniz, dan sifat-sifat lainnya semuanya
	dapat diperiksa pada $(X,d)$. Dengan cara yang sama terlihat bahwa
	$\Hm(E_r,d_r)\simeq E_{r+1}$ merupakan isomorfisme aljabar bergradasi.

	Pernyataan bahwa isomorfisme kanonik
	$E_\infty^{p,q}\simeq\gr^p\Hm^{p+q}(X,d)$ dalam teorema konvergensi klasik
	mempertahankan perkalian juga cukup diverifikasi pada $X$, sehingga
	tidak perlu diuraikan lagi.
\end{proof}

% segment-id: o014.li2.chapter5.multiplicative-structures.r001
\begin{remark}
	\index{sistem Cartan--Eilenberg (Cartan--Eilenberg system)}
	Struktur perkalian merupakan alat yang sangat diandalkan dalam topologi.
	Topologi aljabar tanpa perkalian hampir tidak terbayangkan, atau setidaknya
	akan terasa hambar. Secara historis, ketika pertama kali menemukan barisan
	spektral, Leray sudah mempertimbangkan struktur perkalian dan menyebutnya
	``cincin spektral''. Struktur ini sering sangat menyederhanakan perhitungan
	barisan spektral dalam topologi. Itulah alasan bagian ini memberi pengantar
	singkat mengenai struktur perkalian, meskipun hanya menyentuh permukaannya.

	Ada satu kesulitan yang jelas ketika menerapkan
	Definisi~\ref{def:mult-ss}. Struktur aljabar bergradasi diferensial pada
	semua $E_r$ harus diberikan sekaligus, sebab dari definisinya saja tidak ada
	alasan bahwa perkalian pada satu halaman menginduksi perkalian pada halaman
	berikutnya. Dalam sejumlah keadaan hal ini memang dapat dilakukan, seperti
	dalam Proposisi~\ref{prop:dga-mult-ss}.

	Dalam kasus umum atau kasus yang sulit dihitung secara manual, konstruksi
	barisan spektral sering bermula dari data sederhana, misalnya pasangan eksak
dalam \S\ref{sec:exact-couples}. Karena itu wajar ditanyakan: adakah syarat
sederhana pada pasangan eksak yang menjamin bahwa barisan spektral terkait
mempunyai struktur perkalian?

	Tampaknya jawabannya negatif: diperlukan informasi selain pasangan eksak.
Salah satu konstruksi yang berguna adalah \emph{sistem Cartan--Eilenberg} yang
juga klasik, bersama hasil kali spektralnya; lihat \cite[II.A]{Dou58}. Sebagai
penerapan, barisan spektral Serre--Atiyah--Hirzebruch yang sangat penting dalam
topologi mempunyai struktur perkalian untuk setiap teori kohomologi rampat
yang multiplikatif. Karena materi ini sudah menyimpang dari garis utama, kita
berhenti sampai di sini.
\end{remark}
